% page 33 of the facsimile (image p-032 of the scan)
% block 170.2 x 246.3 mm ; left margin 31.8 mm
\pgc{10.160mm}{0.000mm}{
\l{\cel{56}25}
\l{}
\l{}
\l{}
\l{\sou{\textquotedbl{}ampère}\textquotedbl{} (\sou{elektr.}) Intenseso-grado di elektro-fluo qua povas}
\l{\cel{3}deskompozar 0,001118 gramo di arjento singla-sekunde, per}
\l{\cel{3}elektrolizo, ek aquo-dissolvuro de arjento-nitrato. -}
\l{\cel{3}DEFIRS.}
\l{}
\l{\sou{ampèremetro}.- (\sou{elektro)}.\cel{2}Instrumento por mezurar la intenseso-}
\l{\cel{3}grado di elektro-fluo, e di qua la cifro-plako konoernas le}
\l{\cel{3}\textquotedbl{}ampère\textquotedbl{}. - DEFIRS.}
\l{}
\l{\sou{ampla}. Vorto qua expresas relato inter du kozi de speci diversa,}
\l{\cel{3}maxim-multa-kaze la relato di kontenanto e kontenajo.}
\l{\cel{3}\sou{\textquotedbl{}Amplan\textquotedbl{}}\cel{2}qualifikesas :\cel{1}\sou{mantelo}\cel{1}qua fitas tre laxe olua}
\l{\cel{3}+weranto\textquotesingle{};\cel{1}\sou{chambro}\cel{1}od\cel{1}\sou{agro}\cel{1}ube granda asemblitaro esos prezenta:}
\l{\cel{3}\sou{provizuro}\cel{1}qua satisfacas la demandi. Stato qua havas granda}
\l{\cel{3}mar-komerco mustas havar navaro \textquotedbl{}ampla\textquotedbl{}. - EFIS.}
\l{}
\l{\sou{amplitudo.}\cel{2}Grandeso-grado di arko. - DEFIRS.}
\l{}
\l{\sou{ampulo.}\cel{1}I. Mikra fialo, pasable sfera, qua finas pinte o quan}
\l{\cel{3}onu klozas per weldo-lampo pos enduktir en olu liquido. -}
\l{\cel{3}II. Vitra vazo qua kontenas la fili di elektro-lampo.\cel{2}III}
\l{\cel{3}veziketo produktata da la su-akumulo di seraji sub la epidermo.}
\l{\cel{3}- EFIS.}
\l{}
\l{\sou{amputar}.\cel{2}(\sou{trans., de}). De-sekar per tranchilo, per skalpelo}
\l{\cel{3}(membro, parto de membro, organo, e c.) - (\sou{anke metaf.})}
\l{\cel{3}- DEFIRS.}
\l{}
\l{\sou{amuleto}.\cel{2}Objekto a qua onu atribuas supersticoze kuraciveso.-}
\l{\cel{3}- DEFIRS.}
\l{}
\l{\sou{amuzar}. (\sou{trans}.) Okupar per kozi qui disipas tempo. - DEF.}
\l{}
\l{\sou{-an -}\cel{1}. Sufixo qua indikas individuo qua apartenas a klaso}
\l{\cel{3}(urbo, lando, ensemblo) - homo qua esas \textquotedbl{}membro di\textquotedbl{}; plu}
\l{\cel{3}precize : qua apartenas ad ula domeno (lando, socio, e c.).}
\l{\cel{3}Konseque, lu ne uzesas nur pri homi, ma anke pri kozi.}
\l{}
\l{\sou{an}. Prepoziciono qua indikas relato di kontigueso o di apogo,}
\l{\cel{3}tale ke la kozo kontaktas o preske. - D.}
\l{}
\l{\sou{anabaptismo}. (\sou{religio kristana})\cel{2}Herezio qua konsistas ye negar}
\l{\cel{3}la efiko di la bapto kande olca eventas ante la racion-evo,}
\l{\cel{3}e ye baptar itere la adulti. - DEFIRS.}
\l{}
\l{\sou{anabaptisto}. Adheranto di anabaptismo. - DEFIRS.}
\l{}
\l{\sou{anabioso}.\cel{2}(\sou{biol}.) Itero di la vivo videbla, pos longa}
\l{\cel{3}periodo de vivo latenta. - DEFIS.}
\l{}
\l{\sou{anado}.\cel{2}(\sou{zool.})\cel{2}Ucelo palmipeda, di la familio \textquotedbl{} lamelo-}
\l{\cel{3}rostri\textquotedbl{}, qua vivas o ne-amansite od amansite. -L.\cel{1}\sou{anas},}
\l{\cel{3}\sou{anatis}. - FISL.}
}
